\documentclass[12pt,a4paper]{exam}
\usepackage[utf8]{inputenc}
\usepackage{amsmath,amssymb,amsfonts}
\usepackage{geometry}
\usepackage{enumitem}
\usepackage{graphicx}
\usepackage{hyperref}
\usepackage{xcolor}
\geometry{a4paper, top=2cm, bottom=2cm, left=2.5cm, right=2.5cm}
\hypersetup{colorlinks=true, linkcolor=blue, urlcolor=blue}
\renewcommand{\thequestion}{\textbf{\arabic{question}}}
\renewcommand{\questionlabel}{\thequestion.}
\pointname{ marks}
\pointsdroppedatright
\bracketedpoints
\setlength{\parindent}{0pt}
\setlength{\emergencystretch}{3em}
\allowdisplaybreaks
\newsavebox{\schmathbox}
\newcommand{\fitmath}[1]{%
  \sbox{\schmathbox}{$\displaystyle #1$}%
  \ifdim\wd\schmathbox>\linewidth
    \resizebox{\linewidth}{!}{\usebox{\schmathbox}}%
  \else
    \usebox{\schmathbox}%
  \fi}

\begin{document}

\thispagestyle{empty}
\begin{center}
  \textbf{\LARGE Diag} \\[0.3cm]
  \textbf{Time Allowed: 3 Hours} \hfill \textbf{Maximum Marks: 80} \\[0.3cm]
  \rule{\textwidth}{0.5pt}
\end{center}

\vspace{0.3cm}

\noindent\textbf{General Instructions:}
\begin{enumerate}
  \item {\normalfont\normalcolor This question paper contains 40 questions. All questions are compulsory.}
  \item {\normalfont\normalcolor Questions are grouped into sections A through E.}
  \item {\normalfont\normalcolor Use of calculators is NOT permitted.}
\end{enumerate}
\bigskip
\noindent\textbf{\large Section A: Multiple Choice \& Assertion-Reason (1 mark each)}

\begin{questions}
\question[1] The distance of the point $P(3, 4)$ from the origin is:
\begin{enumerate}[label=(\Alph*)]
  \item (A) 3
  \item (B) 4
  \item (C) 5
  \item (D) 7
\end{enumerate}
\vspace{0.15cm}

\question[1] The coordinates of the midpoint of the line segment joining the points $A(2, 4)$ and $B(4, 6)$ are:
\begin{enumerate}[label=(\Alph*)]
  \item (A) (2, 3)
  \item (B) (3, 5)
  \item (C) (4, 4)
  \item (D) (6, 10)
\end{enumerate}
\vspace{0.15cm}

\question[1] The abscissa of the point $(5, 8)$ is:
\begin{enumerate}[label=(\Alph*)]
  \item (A) 5
  \item (B) 8
  \item (C) 3
  \item (D) 13
\end{enumerate}
\vspace{0.15cm}

\question[1] If the distance between points $(2, -2)$ and $(-1, x)$ is 5, then the value of $x$ is not asked, but what is the distance formula expression?
\begin{enumerate}[label=(\Alph*)]
  \item (A) $\sqrt{(2-1)^2 + (-2-x)^2}$
  \item (B) $\sqrt{(2+1)^2 + (-2+x)^2}$
  \item (C) $\sqrt{(1-2)^2 + (x-2)^2}$
  \item (D) $\sqrt{(-1-2)^2 + (x+2)^2}$
\end{enumerate}
\vspace{0.15cm}

\question[1] The point which lies on the y-axis is:
\begin{enumerate}[label=(\Alph*)]
  \item (A) (2, 0)
  \item (B) (1, 1)
  \item (C) (0, 5)
  \item (D) (5, 5)
\end{enumerate}
\vspace{0.15cm}

\question[1] Assertion (A): The point $(0, 5)$ lies on the y-axis. Reason (R): For any point on the y-axis, the x-coordinate is zero.
\begin{enumerate}[label=(\Alph*)]
  \item (A) Both Assertion (A) and Reason (R) are true and Reason (R) is the correct explanation of Assertion (A)
  \item (B) Both Assertion (A) and Reason (R) are true but Reason (R) is NOT the correct explanation of Assertion (A)
  \item (C) Assertion (A) is true but Reason (R) is false
  \item (D) Assertion (A) is false but Reason (R) is true
\end{enumerate}
\vspace{0.15cm}

\question[1] Assertion (A): The distance of the point $P(3, 4)$ from the origin $(0, 0)$ is $5$ units. Reason (R): The distance of a point $(x, y)$ from the origin is calculated by $\sqrt{x^2 + y^2}$.
\begin{enumerate}[label=(\Alph*)]
  \item (A) Both Assertion (A) and Reason (R) are true and Reason (R) is the correct explanation of Assertion (A)
  \item (B) Both Assertion (A) and Reason (R) are true but Reason (R) is NOT the correct explanation of Assertion (A)
  \item (C) Assertion (A) is true but Reason (R) is false
  \item (D) Assertion (A) is false but Reason (R) is true
\end{enumerate}
\vspace{0.15cm}

\question[1] Assertion (A): The midpoint of the line segment joining $(2, 4)$ and $(6, 8)$ is $(4, 6)$. Reason (R): The midpoint formula is $(\frac{x_1+x_2}{2}, \frac{y_1+y_2}{2})$.
\begin{enumerate}[label=(\Alph*)]
  \item (A) Both Assertion (A) and Reason (R) are true and Reason (R) is the correct explanation of Assertion (A)
  \item (B) Both Assertion (A) and Reason (R) are true but Reason (R) is NOT the correct explanation of Assertion (A)
  \item (C) Assertion (A) is true but Reason (R) is false
  \item (D) Assertion (A) is false but Reason (R) is true
\end{enumerate}
\vspace{0.15cm}

\question[1] Assertion (A): The points $(1, 1), (2, 2)$ and $(3, 3)$ are collinear. Reason (R): Three points are collinear if the sum of distances between any two pairs is equal to the distance between the third pair.
\begin{enumerate}[label=(\Alph*)]
  \item (A) Both Assertion (A) and Reason (R) are true and Reason (R) is the correct explanation of Assertion (A)
  \item (B) Both Assertion (A) and Reason (R) are true but Reason (R) is NOT the correct explanation of Assertion (A)
  \item (C) Assertion (A) is true but Reason (R) is false
  \item (D) Assertion (A) is false but Reason (R) is true
\end{enumerate}
\vspace{0.15cm}

\question[1] Assertion (A): The area of a triangle with vertices $(0, 0), (4, 0),$ and $(0, 3)$ is $6$ sq units. Reason (R): Area of a triangle with coordinates $(0, 0), (x_1, y_1), (x_2, y_2)$ is $\frac{1}{2}|x_1y_2 - x_2y_1|$.
\begin{enumerate}[label=(\Alph*)]
  \item (A) Both Assertion (A) and Reason (R) are true and Reason (R) is the correct explanation of Assertion (A)
  \item (B) Both Assertion (A) and Reason (R) are true but Reason (R) is NOT the correct explanation of Assertion (A)
  \item (C) Assertion (A) is true but Reason (R) is false
  \item (D) Assertion (A) is false but Reason (R) is true
\end{enumerate}
\vspace{0.15cm}

\question[1] Assertion (A): If $\sin \theta = \frac{1}{2}$, then $\theta = 30^\circ$. Reason (R): The value of $\sin \theta$ increases as $\theta$ increases from $0^\circ$ to $90^\circ$.
\begin{enumerate}[label=(\Alph*)]
  \item (A) Both Assertion (A) and Reason (R) are true and Reason (R) is the correct explanation of Assertion (A)
  \item (B) Both Assertion (A) and Reason (R) are true but Reason (R) is NOT the correct explanation of Assertion (A)
  \item (C) Assertion (A) is true but Reason (R) is false
  \item (D) Assertion (A) is false but Reason (R) is true
\end{enumerate}
\vspace{0.15cm}

\question[1] Assertion (A): The value of $\cos \theta$ can never be $1.5$. Reason (R): The range of $\cos \theta$ is $[-1, 1]$.
\begin{enumerate}[label=(\Alph*)]
  \item (A) Both Assertion (A) and Reason (R) are true and Reason (R) is the correct explanation of Assertion (A)
  \item (B) Both Assertion (A) and Reason (R) are true but Reason (R) is NOT the correct explanation of Assertion (A)
  \item (C) Assertion (A) is true but Reason (R) is false
  \item (D) Assertion (A) is false but Reason (R) is true
\end{enumerate}
\vspace{0.15cm}

\question[1] Assertion (A): $\sin^2 \theta + \cos^2 \theta = 1$ for all $\theta$. Reason (R): $\tan^2 \theta - \sec^2 \theta = 1$.
\begin{enumerate}[label=(\Alph*)]
  \item (A) Both Assertion (A) and Reason (R) are true and Reason (R) is the correct explanation of Assertion (A)
  \item (B) Both Assertion (A) and Reason (R) are true but Reason (R) is NOT the correct explanation of Assertion (A)
  \item (C) Assertion (A) is true but Reason (R) is false
  \item (D) Assertion (A) is false but Reason (R) is true
\end{enumerate}
\vspace{0.15cm}

\question[1] Assertion (A): If $\tan A = \cot A$, then $A = 45^\circ$. Reason (R): $\tan A = \cot(90^\circ - A)$.
\begin{enumerate}[label=(\Alph*)]
  \item (A) Both Assertion (A) and Reason (R) are true and Reason (R) is the correct explanation of Assertion (A)
  \item (B) Both Assertion (A) and Reason (R) are true but Reason (R) is NOT the correct explanation of Assertion (A)
  \item (C) Assertion (A) is true but Reason (R) is false
  \item (D) Assertion (A) is false but Reason (R) is true
\end{enumerate}
\vspace{0.15cm}

\question[1] Assertion (A): The value of $\sin 0^\circ$ is $0$. Reason (R): The value of $\cos 0^\circ$ is $0$.
\begin{enumerate}[label=(\Alph*)]
  \item (A) Both Assertion (A) and Reason (R) are true and Reason (R) is the correct explanation of Assertion (A)
  \item (B) Both Assertion (A) and Reason (R) are true but Reason (R) is NOT the correct explanation of Assertion (A)
  \item (C) Assertion (A) is true but Reason (R) is false
  \item (D) Assertion (A) is false but Reason (R) is true
\end{enumerate}
\vspace{0.15cm}

\end{questions}
\bigskip
\noindent\textbf{\large Section B: Very Short Answer (2 marks each)}

\begin{questions}
\question[2] In $\triangle PQR$, $S$ and $T$ are points on sides $PQ$ and $PR$ respectively such that $\angle PST = \angle PRQ$. If $PS = 2\text{ cm}$, $SQ = 3\text{ cm}$, and $PT = 2.5\text{ cm}$, find $TR$.
\vspace{0.15cm}

\question[2] Diagonals of a trapezium $\text{ABCD}$ with $AB \parallel DC$ intersect each other at the point $O$. If $AB = 2CD$, find the ratio of the areas of triangles $\triangle AOB$ and $\triangle COD$.
\vspace{0.15cm}

\end{questions}
\bigskip
\noindent\textbf{\large Section C: Short Answer (3 marks each)}

\begin{questions}
\question[3] In $\triangle ABC$, $DE \parallel BC$ such that $AD = (4x - 3)$ cm, $AE = (8x - 7)$ cm, $BD = (3x - 1)$ cm, and $EC = (5x - 3)$ cm. Find the value of $x$.
\vspace{0.15cm}

\question[3] Two poles of heights $6$ m and $11$ m stand on a plane ground. If the distance between their feet is $12$ m, find the distance between their tops.
\vspace{0.15cm}

\question[3] In $\triangle PQR$, $S$ and $T$ are points on sides $PQ$ and $PR$ respectively such that $\angle PST = \angle PRQ$. Show that $\triangle PQR \sim \triangle PTS$.
\vspace{0.15cm}

\question[3] The areas of two similar triangles are $121$ cm$^2$ and $64$ cm$^2$ respectively. If the median of the first triangle is $12.1$ cm, find the corresponding median of the other triangle.
\vspace{0.15cm}

\question[3] In $\triangle ABC$, $AD \perp BC$ such that $AD^2 = BD \cdot CD$. Prove that $\angle BAC = 90^\circ$.
\vspace{0.15cm}

\question[3] Find the ratio in which the line segment joining the points $A(-3, 10)$ and $B(6, -8)$ is divided by the point $P(-1, 6)$.
\vspace{0.15cm}

\question[3] If the points $A(6, 1)$, $B(8, 2)$, $C(9, 4)$, and $D(p, 3)$ are the vertices of a parallelogram taken in order, find the value of $p$.
\vspace{0.15cm}

\question[3] Find the distance between the points $P(a+b, b-a)$ and $Q(a-b, a+b)$.
\vspace{0.15cm}

\question[3] Find the value of $y$ for which the distance between the points $P(2, -3)$ and $Q(10, y)$ is 10 units.
\vspace{0.15cm}

\question[3] Find the coordinates of a point $A$, where $AB$ is the diameter of a circle whose centre is $(2, -3)$ and $B$ is $(1, 4)$.
\vspace{0.15cm}

\end{questions}
\bigskip
\noindent\textbf{\large Section D: Long Answer (4-5 marks each)}

\begin{questions}
\question[5] State and prove the Basic Proportionality Theorem (Thales Theorem).
\vspace{0.15cm}

\question[5] Prove that the ratio of the areas of two similar triangles is equal to the square of the ratio of their corresponding sides.
\vspace{0.15cm}

\question[5] In $\triangle ABC$, $AD \perp BC$ such that $AD^2 = BD \cdot CD$. Prove that $\triangle ABC$ is a right-angled triangle at $A$.
\vspace{0.15cm}

\question[5] A vertical stick $2$ m long casts a shadow $1$ m long on the ground. At the same time, a tower casts a shadow $5$ m long. Find the height of the tower using the properties of similar triangles.
\vspace{0.15cm}

\question[5] Find the ratio in which the point $P(x, 2)$ divides the line segment joining the points $A(12, 5)$ and $B(4, -3)$. Also, find the value of $x$.
\vspace{0.15cm}

\question[5] If the points $A(6, 1)$, $B(8, 2)$, $C(9, 4)$, and $D(p, 3)$ are the vertices of a parallelogram taken in order, find the value of $p$.
\vspace{0.15cm}

\question[5] Find the area of a quadrilateral whose vertices, taken in order, are $(-4, -2)$, $(-3, -5)$, $(3, -2)$, and $(2, 3)$.
\vspace{0.15cm}

\question[5] Find the coordinates of the points which divide the line segment joining $A(-2, 2)$ and $B(2, 8)$ into four equal parts.
\vspace{0.15cm}

\question[1] If $\sin \theta = \frac{3}{5}$, then the value of $\cos \theta$ is:
\vspace{0.15cm}

\question[1] The value of $\tan 45^\circ + \cot 45^\circ$ is:
\vspace{0.15cm}

\question[1] If $3\tan \theta = 3$, then $\theta$ is equal to:
\vspace{0.15cm}

\question[1] The value of $\frac{\sin 18^\circ}{\cos 72^\circ}$ is:
\vspace{0.15cm}

\question[1] Which of the following is equal to $\sec^2 \theta - 1$?
\vspace{0.15cm}

\end{questions}
\bigskip

\end{document}
